\documentclass[10pt,letterpaper]{article}

\usepackage[T1]{fontenc}
\usepackage[utf8]{inputenc}
\usepackage[spanish,es-noquoting]{babel}
\usepackage{lmodern}
\usepackage[margin=1.9cm,top=1.5cm,bottom=1.6cm]{geometry}
\usepackage{booktabs}
\usepackage{array}
\usepackage{xcolor}

\definecolor{acc}{HTML}{1F4E79}
\definecolor{gris}{HTML}{5B6672}

\setlength{\parskip}{2pt}
\setlength{\topsep}{2pt}
\setlength{\partopsep}{0pt}
\setlength{\parindent}{0pt}
\renewcommand{\arraystretch}{1.05}

% Secciones y pie de pagina sin paquetes externos: solo LaTeX base.
\makeatletter
\renewcommand\section{\@startsection{section}{1}{\z@}%
  {7pt \@plus 2pt \@minus 1pt}{3pt}%
  {\normalfont\normalsize\bfseries\color{acc}}}
\renewcommand\@seccntformat[1]{\csname the#1\endcsname.\quad}
\renewcommand{\ps@plain}{%
  \renewcommand{\@oddhead}{}\renewcommand{\@evenhead}{}%
  \renewcommand{\@oddfoot}{\footnotesize\color{gris}%
    TI-4601 $\cdot$ Laboratorio 1 $\cdot$ Cl\'uster CockroachDB, localidad y qu\'orum
    \hfill \thepage}%
  \renewcommand{\@evenfoot}{\@oddfoot}}
\makeatother
\pagestyle{plain}

\newcommand{\nodo}[3]{%
  \fbox{\parbox[t][1.35cm][c]{0.29\textwidth}{\centering
    \textbf{#1}\\[2pt]
    {\ttfamily\footnotesize #2}\\[2pt]
    {\footnotesize\color{gris}\itshape #3}}}}

\begin{document}

{\LARGE\bfseries\color{acc} Laboratorio 1 --- Clúster CockroachDB, localidad y quórum}\\[2pt]
{\footnotesize\color{gris} TI-4601 Bases de Datos Avanzadas $\cdot$ Grupo 3 $\cdot$ Fabricio Mena $\cdot$ Christopher 	Rodriguez $\cdot$ Alisson Redondo}

\vspace{4pt}\hrule\vspace{6pt}

\section{Nodos, localidades, \texttt{-{}-join} y \texttt{-{}-locality}}

\begin{center}
\nodo{crdb-1}{region=cr-sj}{SQL 26257 / UI 8080}\hfill
\nodo{crdb-2}{region=cr-limon}{sin puertos publicados}\hfill
\nodo{crdb-3}{region=us-east}{sin puertos publicados}
\end{center}
\vspace{-4pt}
\begin{center}
{\footnotesize\color{gris} Los tres con \texttt{-{}-join=crdb-1,crdb-2,crdb-3} y \texttt{zone=a};
volúmenes \texttt{crdb1}, \texttt{crdb2}, \texttt{crdb3}. Una sola computadora: las tres regiones son
localidades lógicas, no una WAN.}
\end{center}

\textbf{\texttt{-{}-join}} es la lista con la que los tres procesos se encuentran y forman un solo
clúster; no declara regiones. Tras \texttt{make lab1-up} los tres nodos aparecen vivos en
\texttt{node-status-initial.txt}, pero \texttt{SHOW REGIONS FROM DATABASE ti4601} no devuelve nada
hasta ejecutar el Paso 3: formar el clúster y configurar la base multi-región son cosas distintas.

\textbf{\texttt{-{}-locality}} es una etiqueta lógica, el par \texttt{region=...,zone=...}.
CockroachDB la usa para colocar réplicas y para resolver \texttt{REGIONAL BY ROW} y \texttt{GLOBAL},
pero el nombre no crea distancia: etiqueta no es WAN física. Borrar los volúmenes costaría el
almacenamiento estable que Raft supone --- término actual, voto del término y log replicado ---
además de regiones, tablas y datos.

\section{Regiones, esquema y semilla}

{\footnotesize
\begin{verbatim}
ALTER DATABASE ti4601 PRIMARY REGION "cr-sj";   ALTER DATABASE ti4601 ADD REGION "cr-limon";
ALTER DATABASE ti4601 ADD REGION "us-east";     SHOW REGIONS FROM DATABASE ti4601;
\i labs/lab1-cluster/schema.sql    \i labs/lab1-cluster/seed.sql    \i labs/lab1-cluster/raft_probe.sql
\end{verbatim}}
\vspace{2pt}

Cockroach generó \texttt{LOCALITY GLOBAL} para \texttt{catalogo\_producto} y
\texttt{LOCALITY REGIONAL BY ROW AS region} para \texttt{pedido}. \texttt{make lab1-check} terminó
en \texttt{6/6} (\texttt{config-check.txt}).

\textbf{Decisión 1: \texttt{region} va primero en la clave primaria de \texttt{pedido}.}
\texttt{REGIONAL BY ROW AS region} reparte las filas por esa columna. Al ir primero en
\texttt{PRIMARY KEY (region, pedido\_id)}, las filas de una región quedan juntas y esa región puede
ser su hogar. El costo es que una consulta que no filtre por \texttt{region} puede producir
\emph{scatter}.

\textbf{Decisión 2: \texttt{catalogo\_producto} es \texttt{GLOBAL}, no \texttt{REGIONAL BY ROW}.}
Es el catálogo pequeño que se quiere leer desde las tres regiones. \texttt{GLOBAL} da esa lectura
en todas ellas; el precio es que la lectura global tiene costo de escritura. Con
\texttt{REGIONAL BY ROW} habría que elegir una sola región hogar.

\section{Réplica Raft frente a fragmentación y asignación}

\begin{center}\footnotesize
\begin{tabular}{@{}p{0.30\textwidth}p{0.32\textwidth}p{0.31\textwidth}@{}}
\toprule
\textbf{Diseño de fragmentación} & \textbf{En CockroachDB} & \textbf{No confundir} \\
\midrule
Fragmentación horizontal primaria por región & \texttt{pedido}: \texttt{REGIONAL BY ROW AS region} & hogar de fila $\neq$ consenso \\
Catálogo pequeño replicado & \texttt{catalogo\_producto}: \texttt{GLOBAL} & lectura global tiene costo de write \\
Sitio / región & \texttt{-{}-locality=region=...} & etiqueta $\neq$ WAN física \\
Réplica selectiva & política de colocación & réplica Raft copia un rango \\
\bottomrule
\end{tabular}
\end{center}

Fragmentar y asignar son decisiones de diseño sobre relaciones: qué filas existen y en qué sitio
viven. Una réplica Raft es otra cosa: una copia de un \textbf{rango}, mantenida por un grupo de
consenso para sobrevivir a la caída de un nodo. Son fenómenos distintos en el mismo clúster.
\texttt{ti4601\_raft.public.raft\_probe} no usa localidad multirregional y aun así tiene tres
réplicas votantes; por eso la disponibilidad del §5 no debe atribuirse a \texttt{REGIONAL BY ROW}.

\textbf{Varios rangos implican varios grupos Raft.} El motor parte el espacio de claves en rangos y
cada rango tiene sus propias réplicas, su líder Raft y su leaseholder: es una máquina de estados
replicada independiente. En \texttt{cluster-inspect.txt}, \texttt{raft\_probe} ocupa el rango
\texttt{76} con sus votantes \texttt{\{1,3,2\}}; las filas de \texttt{pedido} viven en otros rangos,
con sus propios grupos. El paper explica un log; el motor ejecuta muchos (Multi-Raft), y una
consulta puede tocar varios rangos.

\textbf{Límite de este clúster.} \texttt{REGIONAL BY ROW} espera varios nodos por región y aquí hay
uno, así que algunos rangos de \texttt{pedido} pueden aparecer subreplicados: la evidencia demuestra
región hogar y localidad configurada, no residencia estricta.

\section{Latencia: p50 y p99 de los cuatro casos}

Gateway fijo \texttt{crdb-1}, warm-up de 5 descartado por caso y \texttt{n=50} muestras por caso
(200 en total; \texttt{latency.csv} tiene 201 líneas). El reloj es
\texttt{time.perf\_counter\_ns()} y el p99 usa \emph{nearest rank}.

\begin{center}\footnotesize
\begin{tabular}{@{}llcrrr@{}}
\toprule
\textbf{operation} & \textbf{locality} & \textbf{home\_region} & \textbf{n} & \textbf{p50\_ms} & \textbf{p99\_ms} \\
\midrule
read  & local  & cr-sj    & 50 & 0.690 & 1.934 \\
read  & remote & cr-limon & 50 & 1.460 & 1.698 \\
write & local  & cr-sj    & 50 & 7.067 & 12.713 \\
write & remote & cr-limon & 50 & 6.072 & 11.792 \\
\bottomrule
\end{tabular}
\end{center}

Lo medido se reporta tal cual. En escrituras el caso remoto salió incluso por debajo del local:
resultado válido, porque Docker no añade distancia física y la razón local/remoto puede quedar cerca
de 1 o invertirse por ruido. La diferencia clara está entre lectura y escritura, no entre local y
remoto.

\textbf{Declaración obligatoria.} (i) Las tres regiones son lógicas. (ii) La máquina, los recursos
disponibles y el número de ejecuciones se mantuvieron constantes. (iii) Inferir el costo de una WAN
requeriría latencia controlada, ausente hoy.

\section{Escenario A: falla de un nodo, RTO y RPO}

\textbf{Precondición} (\texttt{chaos-a-before.txt}). El rango \texttt{76} de \texttt{raft\_probe}
tenía \texttt{voting\_replicas \{1,3,2\}}, una por nodo. El lease se movió al nodo que se iba a
detener con \texttt{ALTER RANGE RELOCATE LEASE TO 3} (el \texttt{store\_id} de
\texttt{region=cr-limon}), y quedó con \texttt{lease\_holder = 3}. Ese \texttt{node 3} es el
contenedor \texttt{ti4601-crdb-2}: se detuvo el leaseholder, no un follower. El valor confirmado antes
de la falla era \texttt{version = 0}.

\begin{center}\footnotesize
\begin{tabular}{@{}lll@{}}
\toprule
\textbf{Instante (UTC)} & \textbf{Evento} & \textbf{Fuente} \\
\midrule
17:37:40.543 -- 17:37:47.039 & 14 escrituras \texttt{before-stop ok} & \texttt{chaos-a.csv} \\
\textbf{17:37:47.5218} & \texttt{docker stop -{}-timeout 0 ti4601-crdb-2} & \texttt{chaos-a-stop.epoch} \\
17:37:49.539 & \texttt{after-stop error} 2007.6 ms & \texttt{chaos-a.csv} \\
17:37:51.549 & \texttt{after-stop error} 2009.2 ms & \texttt{chaos-a.csv} \\
\textbf{17:37:51.5594} & \textbf{primer \texttt{after-stop ok}}, 10.235 ms & \texttt{chaos-a.csv} \\
17:37:59.472 & \texttt{crdb-2} vuelve a estar arriba & \texttt{chaos-a-node-status.txt} \\
17:38:10.055 & lectura posterior: \texttt{version = 112} & \texttt{chaos-a-rpo.txt} \\
\bottomrule
\end{tabular}
\end{center}

\textbf{RTO observado.} Se toma el instante de falla de \texttt{chaos-a-stop.epoch} y la columna
\texttt{completed\_epoch} de la primera fila con \texttt{phase=after-stop} y \texttt{status=ok}:

{\footnotesize
\begin{verbatim}
RTO = (1788370671.559400 - 1788370667.521754) x 1000 = 4037.646 ms  ->  4037.6 ms
\end{verbatim}}
\vspace{2pt}

La sonda imprimió \texttt{4037.6 ms}: el mismo cálculo, redondeado a un decimal.

\textbf{Errores transitorios: 2}, ambos \texttt{QueryCanceled} por el \texttt{statement\_timeout}
de 2 s de la sonda. Son intentos sin confirmación, de resultado desconocido; no se contabilizan
como operaciones confirmadas perdidas.

\textbf{RPO observado para operaciones confirmadas: 0.} \texttt{version} solo se incrementa. El
valor confirmado antes de la falla no retrocedió: la lectura posterior a la recuperación devuelve
\texttt{version = 112}, mayor que el valor previo, y \texttt{make lab1-status} muestra los tres
nodos con \texttt{is\_live = true}.

\section{Mayoría de 2 de 3}

Con tres votantes la mayoría es 2, porque dos mayorías cualesquiera comparten al menos un servidor.
Esa intersección es la que hace que un write committed sobreviva: el próximo líder necesita votos de
una mayoría, y esa mayoría contiene a alguien que ya tiene la entrada. Mayoría no significa muchas
copias; significa intersección.

La evidencia del §5 lo muestra: las escrituras volvieron a confirmar a las 17:37:51.56 y
\texttt{crdb-2} no regresó sino hasta las 17:37:59.47. Las 38 escrituras \texttt{after-stop ok} se
confirmaron con dos nodos de tres, no porque el nodo caído volviera.

Con dos nodos caídos queda un proceso activo pero 1 de 3, menos que la mayoría: sigue vivo y acepta
la conexión, pero no puede reunir mayoría para elegir líder ni para avanzar \texttt{commitIndex}, así
que la escritura no debe confirmarse. Proceso activo no es lo mismo que quórum.

\section{Puente al Proyecto 1}

Se reutiliza el método, no los artefactos: Compose y la imagen cliente, las variables \texttt{PG*},
\texttt{SHOW REGIONS}, \texttt{SHOW CREATE} y \texttt{SHOW RANGES}, las definiciones de p50, p99,
RTO y RPO, y la estructura precondición $\rightarrow$ evento $\rightarrow$ medición $\rightarrow$
verificación. Entregar este DDL renombrado no cumple el P1.

Queda pendiente y es trabajo propio del equipo: el dominio del E1 con sus predicados, claves y
localidades, y su propia tabla \texttt{REGIONAL BY ROW} y su candidato a \texttt{GLOBAL}; un script
reproducible de configuración, porque aquí se aplicó a mano en \texttt{psql}; un seed y un generador
de carga propios; un programa que mida operaciones representativas del dominio; una tabla operacional
propia con RF=3 para el E4, distinta de \texttt{raft\_probe}; la orquestación de la falla y su
evidencia; y el contraste del E5 contra un nodo o el Lab 0. Los límites que deben declararse: las
regiones son lógicas, los rangos \texttt{REGIONAL BY ROW} pueden quedar subreplicados con un nodo por
región, y de este Compose no se sigue ninguna afirmación sobre residencia de los datos.

\end{document}
